\documentclass[11pt,reqno,oneside]{amsart}
\pagestyle{plain}
\usepackage[utf8]{inputenc}
\usepackage[T1]{fontenc}
\usepackage{lmodern}
\usepackage{textcomp}
\usepackage{upquote}
\usepackage{amsmath,amssymb,amsthm}
\usepackage{mathtools}
\usepackage{microtype}
\usepackage{enumitem}
\usepackage{needspace}
\usepackage{booktabs}
\usepackage{array}
\usepackage{longtable}
\usepackage[hypcap=false]{caption}
\usepackage[margin=1.1in]{geometry}
\addtolength{\footskip}{15pt}
\usepackage{xcolor}

\definecolor{linkblue}{RGB}{0,0,160}
\usepackage[colorlinks=true,linkcolor=linkblue,citecolor=linkblue,urlcolor=linkblue]{hyperref}

\setlist[itemize]{topsep=2.5pt, partopsep=0pt, itemsep=2pt, parsep=0pt}
\setlist[enumerate]{topsep=2.5pt, partopsep=0pt, itemsep=2pt, parsep=0pt}

\theoremstyle{plain}
\newtheorem{theorem}{Theorem}[section]
\newtheorem{proposition}[theorem]{Proposition}
\newtheorem{lemma}[theorem]{Lemma}
\newtheorem{corollary}[theorem]{Corollary}

\theoremstyle{definition}
\newtheorem{definition}[theorem]{Definition}
\newtheorem{convention}[theorem]{Convention}
\newtheorem{remark}[theorem]{Remark}

\numberwithin{equation}{section}

\newcommand{\N}{\mathbb{N}}
\newcommand{\Z}{\mathbb{Z}}
\newcommand{\Q}{\mathbb{Q}}
\newcommand{\Fp}{\mathbb{Z}/p\mathbb{Z}}
\newcommand{\C}{\mathcal{C}}
\newcommand{\D}{\mathcal{D}}
\newcommand{\reaches}{\rightsquigarrow}

\DeclareMathOperator{\step}{step}
\DeclareMathOperator{\ord}{ord}
\DeclareMathOperator{\Hol}{Hol}
\DeclareMathOperator{\rl}{rl}
\DeclareMathOperator{\llabel}{ll}
\DeclareMathOperator{\mult}{\mu}
\DeclareMathOperator{\Aff}{Aff}

\newcommand{\then}{\mathbin{\fatsemi}}
\makeatletter
\newcommand{\fatsemi}{\mathrel{\raisebox{0.1ex}{\scalebox{0.9}{\ensuremath{\mathbin{;}}}}}}
\makeatother

\title[Every Remainder Everywhere]{Every Remainder Everywhere}
\author[Sio]{Sio}
\date{October 24, 2026}

\begin{document}

\begin{abstract}
We prove that for every prime $p$, every positive integer $w$ not divisible by $3$, every residue class $r \in \mathbb{Z}/p\mathbb{Z}$, and every threshold $B \in \mathbb{N}$, there exists an integer $n > B$ such that $n$ reaches $w$ under iterations of the Collatz map and $n \equiv r \pmod p$. In particular, the set of reaching predecessors of $w$ in any residue class modulo $p$ is infinite. The proof proceeds by analyzing an explicit two-branch affine dynamical system on the $3$-adic integers. We define a family of finite fibres recording the residues modulo $p$ attained by descendants of an adic seed; show that their cardinalities stabilize; and prove that an affine holonomy group acts transitively on the stabilized fibres via a non-trivial translation generated by a group commutator. The entire mathematical development is formalization-faithful, matching a complete verification in Lean~4. Finally, as an application, we establish that single-prime modular obstructions cannot rule out hypothetical Collatz cycle equations.
\end{abstract}

\maketitle

\section*{Standing Assumptions and Conventions}
Throughout this paper, $p$ denotes a fixed prime number. In Sections~\ref{sec:affine} through~\ref{sec:density}, we assume $p \ge 5$. The ring $\Fp$ denotes the finite field of order $p$, and $(n : \Fp)$ denotes the canonical image of $n \in \N$ in $\Fp$. All divisions between natural numbers defining our operators are exact in $\N$ (as established in Lemma~\ref{lem:L_spec}), mirroring the verified formalization in Lean~4. Explicit truncating division $\lfloor \cdot \rfloor$ is used only in Proposition~\ref{prop:transitivity}.

\section{Collatz Dynamics and the Main Result}\label{sec:collatz}

The standard $3n+1$ Collatz operator $\C \colon \N \to \N$ is defined by:
\[
\C(n) = \begin{cases} n/2 & \text{if } n \equiv 0 \pmod 2, \\ 3n+1 & \text{if } n \equiv 1 \pmod 2. \end{cases}
\]
We say $n$ \emph{reaches} $v$ (written $n \reaches v$) if $\C^{(k)}(n) = v$ for some $k \ge 0$.

\begin{theorem}[Main Theorem]\label{thm:main}
Let $p$ be a prime. Let $w \in \N$ with $w > 0$ and $3 \nmid w$, let $r \in \Fp$, and let $B \in \N$. Then there exists $n \in \N$ such that $n > B$, $n \reaches w$, and $(n : \Fp) = r$.
\end{theorem}

The proof of Theorem~\ref{thm:main} splits into two regimes:
\begin{itemize}
    \item \textbf{Primes $p \ge 5$:} We construct an abstract descendant tree $\D(v)$ from a seed $v \in G$ reaching $w$, and establish a density theorem (Theorem~\ref{thm:density}): in every valid cylinder at depth $K$, descendants of $v$ realize \emph{every} residue modulo $p$. Choosing $K$ large enough ensures $n > B$.
    \item \textbf{Small primes $p \in \{2, 3\}$:} Explicit backward paths under pure halvings and division by $3$ resolve all cases directly (Section~\ref{sec:small_primes}).
\end{itemize}

\needspace{8\baselineskip}
\begin{lemma}[Halving and Atomic Inverse Steps]\label{lem:atomic}
\leavevmode
\begin{enumerate}[label=(\arabic*), leftmargin=1.8em]
    \item $\C(2a) = a$, and $\C^{(j)}(2^j m) = m$ for all $j, m \in \N$.
    \item Reachability is transitive: if $a \reaches b$ and $b \reaches c$, then $a \reaches c$.
    \item For all $m, n \in \N$: $\C(m) = n$ if and only if $m = 2n$, or $m$ is odd and $3m + 1 = n$.
\end{enumerate}
\end{lemma}

\begin{proof}
(1) and (2) are immediate. (3) is the definition of $\C$ read backwards: if $m$ is even then $\C(m) = m/2$, so $\C(m) = n$ means $m = 2n$; if $m$ is odd then $\C(m) = 3m+1$.
\end{proof}

Thus a backward step from $n$ has at most two options: \emph{double} ($n \mapsto 2n$, always available), or \emph{undo $3m+1$} (available only when $n = 3m+1$ with $m$ odd). Everything that follows is a bookkeeping of compositions of these two atomic moves.

\begin{corollary}[Multiples of $3$ Admit Only Doubling]\label{cor:three_dead}
If $3 \mid n$ and $\C(m) = n$, then $m = 2n$.
\end{corollary}

\begin{proof}
If $m$ is odd with $3m+1 = n$, then $n \equiv 1 \pmod 3$, contradicting $3 \mid n$. By Lemma~\ref{lem:atomic}(3), $m = 2n$.
\end{proof}

The corollary is the reason the hypothesis $3 \nmid w$ appears in Theorem~\ref{thm:main}, and the reason only odd numbers prime to $3$ can branch.

\section{Good Numbers and the Coordinate \texorpdfstring{$x = (3n+1)/2$}{x = (3n+1)/2}}\label{sec:coords}

\begin{definition}[Good Numbers and the Coordinate]\label{def:good}
An integer $n \in \N$ is \emph{good} if $n$ is odd and $3 \nmid n$. For odd $n$ put
\[
x(n) = \frac{3n+1}{2}.
\]
\end{definition}

\begin{lemma}[Coordinate Step]\label{lem:coord_step}
If $n$ is odd, then $\C(n) = 2x(n)$ and $\C(2x(n)) = x(n)$. In particular $n \reaches x(n)$ in two steps.
\end{lemma}

\begin{proof}
$n$ odd gives $\C(n) = 3n+1 = 2x(n)$, and then $\C(2x(n)) = x(n)$ by Lemma~\ref{lem:atomic}(1).
\end{proof}

The coordinate $x$ fuses the two forced steps $n \mapsto 3n+1 \mapsto (3n+1)/2$ into one. Only the odd, $3$-free numbers are branch points, and the next lemma identifies them in $x$-coordinates.

\begin{definition}[Cylinders and Valid Classes]\label{def:classes}
A natural number $c \in \N$ is called \emph{valid} if $c \equiv 2 \pmod 9$ or $c \equiv 8 \pmod 9$. A \emph{valid class at depth $K$} (for $K \ge 1$) is a valid integer satisfying the boundedness condition $c < 3^{K+1}$. Modulo $3^{K+1}$, the valid classes represent the admissible $3$-adic residue cylinders of our inverse operators.
\end{definition}

\begin{definition}[Seed Set]\label{def:seed_set}
The admissible seed set is
\[
G = \{x \in \N : x \ge 8 \text{ and } (x \equiv 2 \pmod 9 \;\lor\; x \equiv 8 \pmod 9)\}.
\]
\end{definition}

\begin{lemma}[Good Numbers Are the Valid Classes]\label{lem:good_valid}
If $n \ge 5$ is good, then $x(n) \ge 8$ and $x(n)$ is valid. Conversely, if $x \ge 8$ is valid, then $n = (2x-1)/3$ is a good integer with $n \ge 5$ and $x(n) = x$. Hence
\[
G = \{\, x(n) : n \ge 5 \text{ good} \,\}.
\]
\end{lemma}

\begin{proof}
For odd $n$ we have $x(n) \bmod 9$ determined by $n \bmod 18$. The odd residues prime to $3$ modulo $18$ are $1, 5, 7, 11, 13, 17$, and they give $x(n) \equiv 2, 8, 2, 8, 2, 8 \pmod 9$ respectively. The odd residues $3, 9, 15$ (multiples of $3$) give $x(n) \equiv 5, 5, 5 \pmod 9$ and are excluded. The bound $n \ge 5$ gives $x(n) \ge 8$; the only good number below $5$ is $n = 1$, whose coordinate $x(1) = 2$ lies on the trivial cycle $1 \to 4 \to 2 \to 1$ and is excluded by the threshold $x \ge 8$. Conversely, if $x \equiv 2$ or $8 \pmod 9$ then $2x - 1 \equiv 3$ or $15 \pmod{18}$, so $n = (2x-1)/3$ is an integer, odd (as $3n = 2x-1$ is odd), $n \equiv 1$ or $5 \pmod 6$, hence good; $x \ge 8$ gives $n \ge 5$, and $x(n) = x$ by direct computation.
\end{proof}

\begin{remark}[On the Projection]\label{rem:projection}
The coordinate $x$ is a bookkeeping choice, not an extra hypothesis: Lemma~\ref{lem:coord_step} shows each original-coordinate node $n$ reaches its $x$-coordinate in two Collatz steps. Its benefit is that the right moves become pure multiplications ($4x$ or $16x$), so the descendant tree is a clean binary tree on a single set of residue classes. Proposition~\ref{prop:branch_derivation}(3) below shows that the original-coordinate rules ($4n+1$, $16n+5$, and the mod-$9$ left rule) are carried to $R$ and $L$ exactly.
\end{remark}

\section{The Branch Operators as Composite Inverse Steps}\label{sec:branches}

\begin{definition}[Branch Operators and Multipliers]\label{def:operators}
For $x \in \N$, define the right and left inverse branches:
\[
R(x) = \begin{cases} 
4x & \text{if } x \equiv 2 \pmod 9, \\ 
16x & \text{if } x \not\equiv 2 \pmod 9,
\end{cases}
\qquad\text{and}\qquad
L(x) = \frac{\mult(x)(2x-1)}{6},
\]
where the multiplier $\mult(x) \in \{2, 4, 8, 16\}$ is determined by $x \bmod 27$:
\[
\mult(x) = \begin{cases} 
4 & \text{if } x \equiv 2 \text{ or } 20 \pmod{27}, \\
8 & \text{if } x \equiv 8 \pmod{27}, \\
16 & \text{if } x \equiv 11 \pmod{27}, \\
2 & \text{if } x \equiv 17 \text{ or } 26 \pmod{27}.
\end{cases}
\]
\end{definition}

\begin{definition}[Abstract Descendant Tree]\label{def:seed_tree}
For any $v \in G$, the descendant set $\D(v)$ is the smallest subset of $\N$ containing $v$ and closed under $R$ and $L$.
\end{definition}

\begin{proposition}[Derivation of the Branches]\label{prop:branch_derivation}
Let $x \in G$.
\begin{enumerate}[label=(\arabic*), leftmargin=1.8em]
    \item \emph{(Right branch.)} $R(x)$ is the first valid element of the doubling chain $2x, 4x, 8x, 16x$. Explicitly, if $x \equiv 2 \pmod 9$ then $2x$ is not valid and $R(x) = 4x$; if $x \equiv 8 \pmod 9$ then $2x, 4x, 8x$ are not valid and $R(x) = 16x$.
    \item \emph{(Left branch.)} Let $m = (2x-1)/3$. Then $m$ is odd, $\C(m) = 2x$ (the backward $3m+1$ step), and $L(x)$ is the first valid element of $m, 2m, 4m, 8m$.
    \item \emph{(Original coordinates.)} Let $n$ be good and define
    \[
    R_0(n) = \begin{cases} 4n+1 & n \equiv 1 \pmod 3, \\ 16n+5 & n \equiv 2 \pmod 3,\end{cases}
    \qquad
    L_0(n) = \frac{\mult_0(n)\, n - 1}{3},
    \]
    where $\mult_0(n) = 4, 2, 4, 8, 16, 2$ for $n \equiv 1, 2, 4, 5, 7, 8 \pmod 9$. Then $R(x(n)) = x(R_0(n))$ and $L(x(n)) = x(L_0(n))$.
\end{enumerate}
\end{proposition}

\begin{proof}
(1) If $x \equiv 2 \pmod 9$, then $2x, 4x$ have residues $4, 8$ modulo $9$, so the first valid one is $4x$. If $x \equiv 8 \pmod 9$, then $2x, 4x, 8x, 16x$ have residues $7, 5, 1, 2$, so the first valid one is $16x$.

(2) Since $3m = 2x - 1$ is odd, $m$ is odd, and $\C(m) = 3m+1 = 2x$ by Lemma~\ref{lem:atomic}(3). Further $L(x) = \tfrac{\mult(x)}{2}\,m$, by Lemma~\ref{lem:L_spec} below, with $\mult(x)/2 \in \{1, 2, 4, 8\}$. Running through the six classes of $m \bmod 9$ (namely $1, 2, 4, 5, 7, 8$, with $m = n$ in the notation of Lemma~\ref{lem:good_valid}) the first of $m, 2m, 4m, 8m$ that is $\equiv 2, 8 \pmod 9$ is $2m, m, 2m, 4m, 8m, m$ respectively, matching $\mult(x)/2 = 2,1,2,4,8,1$.

(3) For good $n$, $x = x(n)$ satisfies $2x - 1 = 3n$. We have $x \equiv 2 \pmod 9$ iff $3n+1 \equiv 4 \pmod{18}$ iff $n \equiv 1 \pmod 3$. In that case $x(4n+1) = 6n+2 = 4x(n)$; otherwise $x(16n+5) = 24n+8 = 16x(n)$. For $L$: the residues $n \equiv 1, 5, 7, 11, 13, 17 \pmod{18}$ give $x \equiv 2, 8, 11, 17, 20, 26 \pmod{27}$, so $\mult(x) = 4, 8, 16, 2, 4, 2$, which is exactly $\mult_0(n)$ for $n \bmod 9 = 1, 5, 7, 2, 4, 8$. Hence $L(x) = \tfrac{\mult(x)}{2}\,n = \tfrac{\mult_0(n)}{2}\,n$, while $x(L_0(n)) = (3 L_0(n) + 1)/2 = \tfrac{\mult_0(n)}{2}\,n$ (recall $\mult_0(n)$ is even).
\end{proof}

Proposition~\ref{prop:branch_derivation} is the precise sense in which the tree is \emph{derived} from the Collatz map rather than posited: every branch is a composite of the two atomic moves of Lemma~\ref{lem:atomic}(3), cut off at the first valid class.

\begin{lemma}[Branch Trajectories]\label{lem:collatz_elementary}
\leavevmode
\begin{enumerate}[label=(\arabic*), leftmargin=1.8em]
    \item For all $x \in \N$, $R(x) \reaches x$. Specifically, $R(x)$ reaches $x$ in $2$ steps if $x \equiv 2 \pmod 9$, and in $4$ steps otherwise.
    \item If $x \ge 1$ and $x \equiv 2 \pmod 3$, then $L(x) \reaches x$.
\end{enumerate}
\end{lemma}

\begin{proof}
(1) follows from $R(x) = 2^2 x$ or $2^4 x$ and Lemma~\ref{lem:atomic}(1). For (2), by Lemma~\ref{lem:L_spec} below, $6L(x) = \mult(x)(2x-1)$ holds identically. Since $\mult(x) \in \{2, 4, 8, 16\}$, we write $\mult(x) = 2 \cdot 2^i$ with $i \in \{0, 1, 2, 3\}$. Since $x \equiv 2 \pmod 3$, $2x - 1 \equiv 3 \equiv 0 \pmod 3$, so $m = (2x - 1)/3$ is an integer. Thus $6L(x) = 6 \cdot 2^i m$, which gives $L(x) = 2^i m$. Because $3m = 2x - 1$ is odd, $m$ is necessarily odd. Hence $\C(m) = 3m + 1 = 2x$, and $\C(2x) = x$. Therefore, $L(x) = 2^i m$ reaches $x$ in $i + 2$ iterations of $\C$.
\end{proof}

\begin{lemma}[Soundness of the Tree]\label{lem:desc_sound}
If $v \in G$ and $n \in \D(v)$, then $n \reaches v$.
\end{lemma}

\begin{proof}
We argue by induction on the derivation of $n \in \D(v)$. The base case $n = v$ requires $0$ steps. If $n = R(m)$ with $m \in \D(v)$, Lemma~\ref{lem:collatz_elementary}(1) gives $R(m) \reaches m$, and $m \reaches v$ by the induction hypothesis; reachability is transitive. If $n = L(m)$ with $m \in \D(v)$, Lemma~\ref{lem:G_inv} below shows $m \in G$, whence $m \ge 8 \ge 1$ and $m \equiv 2 \pmod 3$. By Lemma~\ref{lem:collatz_elementary}(2), $L(m) \reaches m$, and transitivity yields $L(m) \reaches v$.
\end{proof}

\begin{lemma}[Existence of Admissible Seeds]\label{lem:seed_exist}
\enlargethispage{\baselineskip}
If $w > 0$ and $3 \nmid w$, there exists $v \in G$ such that $v \reaches w$.
\end{lemma}

\begin{proof}
Since $3 \nmid w$, $w \bmod 9 \in \{1, 2, 4, 5, 7, 8\}$. We select $v = 2^k w$ by choosing $k$ according to the following assignment:
\par\vspace{2pt}\noindent\centering
\small
\begin{tabular}{lcccccc}
\toprule
$w \bmod 9$ & $1$ & $2$ & $4$ & $5$ & $7$ & $8$ \\
\midrule
$v$ & $128w$ & $64w$ & $128w$ & $256w$ & $512w$ & $64w$ \\
$k$ & $7$ & $6$ & $7$ & $8$ & $9$ & $6$ \\
$v \bmod 9$ & $2$ & $2$ & $8$ & $2$ & $2$ & $8$ \\
\bottomrule
\end{tabular}
\par\vspace{2pt}\raggedright\noindent
In each case, $v \ge 64 \ge 8$, $v \equiv 2 \text{ or } 8 \pmod 9$, so $v \in G$, and $\C^{(k)}(v) = w$ by Lemma~\ref{lem:atomic}(1).
\end{proof}

\section{Branch Arithmetic and 3-Adic Precision Loss}\label{sec:branch_arith}

\begin{lemma}[Exactness of Left Division]\label{lem:L_spec}
If $x \equiv 2 \pmod 3$, then $6L(x) = \mult(x)(2x-1)$ holds in $\N$.
\end{lemma}

\begin{proof}
By definition, $\mult(x) \in \{2, 4, 8, 16\}$, so $2 \mid \mult(x)$. Since $x \equiv 2 \pmod 3$, $2x - 1 \equiv 2(2) - 1 = 3 \equiv 0 \pmod 3$. Thus $2 \mid \mult(x)$ and $3 \mid (2x - 1)$, which implies $6 \mid \mult(x)(2x-1)$. The division defining $L(x)$ is therefore exact.
\end{proof}

\begin{lemma}[Residue Classification Modulo $27$]\label{lem:mod27}
If $x \equiv 2 \pmod 9$ or $x \equiv 8 \pmod 9$, then
\[
x \bmod 27 \in \{2, 8, 11, 17, 20, 26\}.
\]
\end{lemma}

\begin{proof}
These six residue classes represent the complete preimage of $\{2, 8\}$ under the projection $\Z/27\Z \to \Z/9\Z$.
\end{proof}

\needspace{10\baselineskip}
\begin{lemma}[Invariance of the Domain]\label{lem:G_inv}
Let $x \in G$.
\begin{enumerate}
    \item $R(x) \in G$.
    \item $L(x) \in G$, and $L(x)$ is valid even without the threshold hypothesis $x \ge 8$.
    \item For any $v \in G$ and $n \in \D(v)$, $n \in G$.
\end{enumerate}
\end{lemma}

\begin{proof}
(1) If $x \equiv 2 \pmod 9$, then $R(x) = 4x \equiv 8 \pmod 9$ and $R(x) \ge 4(8) = 32 \ge 8$. If $x \equiv 8 \pmod 9$, then $R(x) = 16x \equiv 7 \cdot 8 = 56 \equiv 2 \pmod 9$ and $R(x) \ge 16(8) \ge 8$.

(2) Since $x \in G$, $x \equiv 2 \pmod 3$, so Lemma~\ref{lem:L_spec} applies. Since $x \equiv a \pmod{27}$ with $a \in \{2, 8, 11, 17, 20, 26\}$, $L(x) \bmod 9$ depends only on $x \bmod 27$. Evaluating $L(a)$ on the representative $a \in \N$ and reducing modulo $9$:
\begin{itemize}[leftmargin=1.5em]
    \item $x \equiv 2 \implies L(2) = 4(3)/6 = 2 \equiv 2$; \quad $x \equiv 8 \implies L(8) = 8(15)/6 = 20 \equiv 2 \pmod 9$;
    \item $x \equiv 11 \implies L(11) = 16(21)/6 = 56 \equiv 2$; \quad $x \equiv 17 \implies L(17) = 2(33)/6 = 11 \equiv 2 \pmod 9$;
    \item $x \equiv 20 \implies L(20) = 4(39)/6 = 26 \equiv 8$; \quad $x \equiv 26 \implies L(26) = 2(51)/6 = 17 \equiv 8 \pmod 9$.
\end{itemize}
In all cases, $L(x) \equiv 2 \text{ or } 8 \pmod 9$. Moreover, for $x \ge 8$, the minimum value is $L(8) = 8(15)/6 = 20 \ge 8$.

(3) Follows by structural induction on $n \in \D(v)$ using (1) and (2).
\end{proof}

\begin{lemma}[Precision Loss under $L$]\label{lem:L_cong}
Let $K \ge 1$. If $x \equiv y \pmod{3^{K+2}}$ with $x \equiv y \equiv 2 \pmod 3$, then $\mult(x) = \mult(y)$ and
\[
L(x) \equiv L(y) \pmod{3^{K+1}}.
\]
\end{lemma}

\begin{proof}
Because $K \ge 1$, $3^3 \mid 3^{K+2}$, so $x \equiv y \pmod{27}$. By Definition~\ref{def:operators}, $\mult(x) = \mult(y)$. Over $\Z$, Lemma~\ref{lem:L_spec} gives:
\[
6(L(x) - L(y)) = \mult(x)(2x - 1) - \mult(x)(2y - 1) = 2\mult(x)(x - y).
\]
Dividing by $2$ in $\Z$ yields $3(L(x) - L(y)) = \mult(x)(x - y)$. Since $3^{K+2} \mid (x - y)$, there exists $q \in \Z$ such that $x - y = 3^{K+2} q$. Substituting and dividing by $3$ gives:
\[
L(x) - L(y) = 3^{K+1} \mult(x) q.
\]
Thus $L(x) \equiv L(y) \pmod{3^{K+1}}$.
\end{proof}

\section{Right-Branch Dynamics on Cylinders}\label{sec:branch_dynamics}

Fix $K \ge 1$. Modulo $3^{K+1}$, the right branch acts as a discrete dynamical system on valid classes. Note that the results in this section are purely $3$-adic and independent of $p$.

\begin{definition}[Cylinder Maps and Step Exponents]\label{def:cylinder_maps}
For a valid class $c < 3^{K+1}$, define $\rho(c) = R(c) \bmod 3^{K+1}$, and for a valid class $e < 3^{K+2}$, define $\lambda(e) = L(e) \bmod 3^{K+1}$. Define the right-step exponent:
\[
\step(c) = \begin{cases} 1 & \text{if } c \equiv 2 \pmod 9, \\ 2 & \text{if } c \not\equiv 2 \pmod 9, \end{cases}
\]
so that $R(c) = c \cdot 4^{\step(c)}$. The $n$-step iterate $\rho^n(c)$ and the cumulative exponent $E(n, c)$ are defined recursively by:
\begin{align*}
\rho^0(c) &= c, & \rho^{n+1}(c) &= \rho^n(\rho(c)), \\
E(0, c) &= 0, & E(n+1, c) &= \step(c) + E(n, \rho(c)).
\end{align*}
\end{definition}

\needspace{6\baselineskip}
\begin{lemma}\label{lem:iter_mult}
Let $c < 3^{K+1}$ be a valid class. For all $n \in \N$:
\[
\rho^n(c) = \bigl(c \cdot 4^{E(n,c)}\bigr) \bmod 3^{K+1}.
\]
\end{lemma}

\begin{proof}
By induction on $n$. The base case $n = 0$ is trivial. For the step, $R(c) = c \cdot 4^{\step(c)}$ and $\rho(c) = R(c) \bmod 3^{K+1}$. Since $\rho(c)$ is valid and bounded, the induction hypothesis gives $\rho^n(\rho(c)) \equiv \rho(c) \cdot 4^{E(n, \rho(c))} \equiv c \cdot 4^{\step(c) + E(n, \rho(c))} \pmod{3^{K+1}}$. Because $\rho^{n+1}(c) < 3^{K+1}$, the remainder coincides with the value.
\end{proof}

\begin{lemma}[Two-Step Tripling]\label{lem:step_pair}
Let $c$ be valid. Then $\step(c) + \step(\rho(c)) = 3$. Consequently, for any valid class $c < 3^{K+1}$ and any $k \in \N$:
\[
E(2k, c) = 3k.
\]
\end{lemma}

\begin{proof}
$R$ alternates the residue modulo $9$: if $c \equiv 2$, then $R(c) \equiv 8$, so $\step(c) + \step(\rho(c)) = 1 + 2 = 3$; if $c \equiv 8$, then $R(c) \equiv 2$, so $2 + 1 = 3$. The second statement follows by induction on $k$.
\end{proof}

\begin{lemma}[Lifting the Exponent for Base $4$ at Prime $3$]\label{lem:lte}
For any $n \ge 1$, $v_3(4^n - 1) = 1 + v_3(n)$. Consequently, for all $K, m \in \N$:
\[
4^m \equiv 1 \pmod{3^{K+1}} \iff 3^K \mid m.
\]
\end{lemma}

\begin{proof}
The valuation identity is the classical LTE lemma for the base $(4, 1)$ at prime $3$:
$v_3(4^n - 1) = v_3(4 - 1) + v_3(n) = 1 + v_3(n)$.
For the equivalence: if $m = 0$, both sides hold trivially. For $m \ge 1$, $4^m \equiv 1 \pmod{3^{K+1}} \iff v_3(4^m - 1) \ge K + 1 \iff 1 + v_3(m) \ge K + 1 \iff 3^K \mid m$.
\end{proof}

\begin{lemma}[Four-Cancellation]\label{lem:four_cancel}
If $a \le b$ and $4^a \equiv 4^b \pmod{3^{K+1}}$, then $a \equiv b \pmod{3^K}$.
\end{lemma}

\begin{proof}
Since $\gcd(4^a, 3^{K+1}) = 1$, we divide by $4^a$ to obtain $4^{b-a} \equiv 1 \pmod{3^{K+1}}$. Lemma~\ref{lem:lte} implies $3^K \mid (b - a)$.
\end{proof}

\begin{lemma}[Exponents via Pigeonhole]\label{lem:exists_exp}
For every valid class $c < 3^{K+1}$, there exists a unique $E \in \{0, 1, \dots, 3^K - 1\}$ such that
\[
c = \bigl(2 \cdot 4^E\bigr) \bmod 3^{K+1}.
\]
Moreover, $c \equiv 2 \pmod 9 \implies E \equiv 0 \pmod 3$, and $c \equiv 8 \pmod 9 \implies E \equiv 1 \pmod 3$.
\end{lemma}

\begin{proof}
Consider the target set $W = \{w < 3^{K+1} : w \equiv 2 \pmod 3\}$. Every $w \in W$ can be uniquely written as $3j + 2$ for $0 \le j < 3^K$, so $|W| = 3^K$. Define $\psi \colon \{0, \dots, 3^K - 1\} \to W$ by $\psi(E) = (2 \cdot 4^E) \bmod 3^{K+1}$. Since $4 \equiv 1 \pmod 3$, $2 \cdot 4^E \equiv 2 \pmod 3$, so $\psi(E) \in W$.

If $\psi(E_1) = \psi(E_2)$, then $2 \cdot 4^{E_1} \equiv 2 \cdot 4^{E_2} \pmod{3^{K+1}}$. Since $\gcd(2, 3^{K+1}) = 1$, $4^{E_1} \equiv 4^{E_2} \pmod{3^{K+1}}$. By Lemma~\ref{lem:four_cancel}, $E_1 \equiv E_2 \pmod{3^K}$. As both are in the range, $E_1 = E_2$. Being an injective map between finite sets of cardinality $3^K$, $\psi$ is surjective. Since valid classes have $c \equiv 2 \pmod 3$, $c \in W$, so $c = \psi(E)$ for some $E$.

For the residue condition: $4^E \bmod 9$ is $1, 4, 7$ according as $E \bmod 3$ is $0, 1, 2$. Therefore, $2 \cdot 4^E \bmod 9$ is $2, 8, 5$. Since $c$ is valid, the residue cannot be $5$, forcing the claimed classification.
\end{proof}

\begin{proposition}[Transitivity of $\rho$]\label{prop:transitivity}
For any two valid classes $c, d < 3^{K+1}$, there exists $n \in \N$ such that $\rho^n(c) = d$.
\end{proposition}

\begin{proof}
By Lemma~\ref{lem:exists_exp}, write $c \equiv 2 \cdot 4^{E_c}$ and $d \equiv 2 \cdot 4^{E_d} \pmod{3^{K+1}}$. Let $q = 3^{K-1} E_c$ (noting $K \ge 1$). Set $y = E_d + 3 \cdot q - E_c = E_d + 3^K E_c - E_c \ge 0$.

We claim there exists $n$ with $E(n, c) = y$. If $y \equiv 0 \pmod 3$, choose $n = 2(y/3)$; by Lemma~\ref{lem:step_pair}, $E(n, c) = 3(y/3) = y$. If $y \equiv 1$ or $2 \pmod 3$, take $n = 2\lfloor y/3 \rfloor + 1$, so $E(n, c) = \step(c) + 3\lfloor y/3 \rfloor$. Reducing $y + E_c = E_d + 3^K E_c$ modulo $3$ gives $y \equiv E_d - E_c$, where $E_c, E_d \in \{0,1\} \pmod 3$ by Lemma~\ref{lem:exists_exp}. If $y \equiv 1$, then $E_c \equiv 0$ and $E_d \equiv 1$, so $c \equiv 2 \pmod 9$ and $\step(c) = 1$. If $y \equiv 2$, then $E_c \equiv 1$ and $E_d \equiv 0$, so $c \equiv 8 \pmod 9$ and $\step(c) = 2$. In both cases $\step(c) \equiv y \pmod 3$, hence $E(n, c) = y$.

By Lemma~\ref{lem:iter_mult}, $\rho^n(c) \equiv c \cdot 4^y \equiv 2 \cdot 4^{E_c + y} \equiv 2 \cdot 4^{E_d + 3^K E_c} \equiv 2 \cdot 4^{E_d} (4^{3^K})^{E_c} \pmod{3^{K+1}}$. By Lemma~\ref{lem:lte}, $4^{3^K} \equiv 1 \pmod{3^{K+1}}$, so $\rho^n(c) \equiv 2 \cdot 4^{E_d} \equiv d \pmod{3^{K+1}}$.
\end{proof}

\needspace{8\baselineskip}
\begin{proposition}[Full Loop at $8$]\label{prop:loop8}
For every $K \ge 1$, there exists $n > 0$ such that $\rho^n(8) = 8$ and $E(n, 8) = 3^K$.
\end{proposition}

\begin{proof}
Write $K = j + 1$. Choose $n = 2 \cdot 3^j$. By Lemma~\ref{lem:step_pair}, $E(n, 8) = 3 \cdot 3^j = 3^{j+1} = 3^K$. Then by Lemma~\ref{lem:iter_mult}, $\rho^n(8) \equiv 8 \cdot 4^{3^K} \equiv 8 \cdot 1 = 8 \pmod{3^{K+1}}$ using Lemma~\ref{lem:lte}. Since $8 < 3^{K+1}$ (as $K \ge 1$), $\rho^n(8) = 8$.
\end{proof}

\section{Affine Labels and Holonomy}\label{sec:affine}

Throughout Sections~\ref{sec:affine} through~\ref{sec:density}, we assume $p \ge 5$. The non-zero elements $2, 3, 4, 6 \in \Fp$ are therefore invertible.

\begin{definition}[Affine Group]\label{def:affine_group}
Define $\Aff(p) = \Fp^\times \times \Fp$. An element $f = (f_1, f_2) \in \Aff(p)$ acts on $\Fp$ by $f(t) = f_1 t + f_2$. Composition in diagrammatic order (first $f$, then $g$) is defined by:
\[
f \then g = (g_1 f_1, \; g_1 f_2 + g_2),
\]
which satisfies $(f \then g)(t) = g(f(t))$. The identity is $(1, 0)$, and the inverse is:
\[
f^{-1} = (f_1^{-1}, \; -f_1^{-1} f_2).
\]
Note that $\Aff(p)$ is strictly non-abelian, and $(f \then g)^{-1} = g^{-1} \then f^{-1}$.
\end{definition}

\begin{definition}[Edge and Walk Labels]\label{def:labels}
Fix depth $K$. The class graph has vertices given by valid classes modulo $3^{K+1}$.
\begin{itemize}
    \item \textbf{Right edges:} For each valid class $c < 3^{K+1}$, an edge $c \to \rho(c)$ with label
    \[
    \rl(c) = \bigl(4^{\step(c)},\; 0\bigr) \in \Aff(p).
    \]
    \item \textbf{Left edges:} For each valid $e < 3^{K+2}$, an edge $(e \bmod 3^{K+1}) \to \lambda(e)$ with label
    \[
    \llabel(e) = \left(\frac{\mult(e)}{3}, \; -\frac{\mult(e)}{6}\right) \in \Aff(p).
    \]
\end{itemize}
A \emph{walk} from $u$ to $w$ is a finite alternating sequence of forward and backward edge traversals. The label of a walk is the composition of its edge labels: forward traversals compose with $f$, backward traversals with $f^{-1}$. The label of the empty walk at $c$ is $(1, 0)$.
\end{definition}

\begin{definition}[Holonomy Group]\label{def:holonomy}
The holonomy group at basepoint $8$ at depth $K$ is:
\[
\Hol(K) = \{f \in \Aff(p) : \text{there exists a closed walk from } 8 \text{ to } 8 \text{ with label } f\}.
\]
By construction, $\Hol(K)$ is a subgroup of $\Aff(p)$.
\end{definition}

\begin{lemma}[Operator Equivariance Modulo $p$]\label{lem:operator_compat}
Let $x \in G$.
\begin{enumerate}
    \item $(R(x) : \Fp) = \rl(x)(x : \Fp)$.
    \item $(L(x) : \Fp) = \llabel(x)(x : \Fp)$.
\end{enumerate}
\end{lemma}

\begin{proof}
(1) Immediate from $R(x) = x \cdot 4^{\step(x)}$.

(2) By Lemma~\ref{lem:L_spec}, $6L(x) = \mult(x)(2x - 1)$ holds in $\N$. Reducing modulo $p$:
\[
6(L(x) : \Fp) = \mult(x)(2(x : \Fp) - 1).
\]
On the other hand,
\[
6 \cdot \llabel(x)(x : \Fp) = 6\left(\frac{\mult(x)}{3}(x : \Fp) - \frac{\mult(x)}{6}\right) = 2\mult(x)(x : \Fp) - \mult(x).
\]
Since $6 \neq 0$ in $\Fp$, we cancel $6$ to obtain the result.
\end{proof}

\needspace{7\baselineskip}
\begin{lemma}[Walks from Right Iterates]\label{lem:walk_of_rIter}
Let $c < 3^{K+1}$ be valid. For any $n \in \N$, there is a forward right-walk from $c$ to $\rho^n(c)$ with label:
\[
\bigl(4^{E(n,c)}, \; 0\bigr) \in \Aff(p).
\]
\end{lemma}

\begin{proof}
Induction on $n$. At $n = 0$, the label is $(4^0, 0) = (1, 0)$. At $n + 1$, prepending the edge $c \to \rho(c)$ with label $(4^{\step(c)}, 0)$ composes to $(4^{E(n, \rho(c))} \cdot 4^{\step(c)}, 0) = (4^{E(n+1, c)}, 0)$.
\end{proof}

\section{Fibre Stabilization and Exact Label Action}\label{sec:fibres}

Fix a seed $v \in G$.

\begin{definition}[Adic Fibres]\label{def:fibres}
For $K \ge 1$ and a valid class $c < 3^{K+1}$, the fibre of $c$ is:
\[
F_K(c) = \{(n : \Fp) \in \Fp : n \in \D(v) \text{ and } n \equiv c \pmod{3^{K+1}}\}.
\]
\end{definition}

\begin{lemma}[Fibre Dynamics]\label{lem:fibre_dynamics}
\begin{enumerate}
    \item For any valid class $c < 3^{K+1}$, $\rl(c)(F_K(c)) \subseteq F_K(\rho(c))$.
    \item For any valid $e < 3^{K+2}$, $\llabel(e)(F_{K+1}(e)) \subseteq F_K(\lambda(e))$.
    \item If $J \le K$ and $c < 3^{K+1}$ is valid, then $F_K(c) \subseteq F_J(c \bmod 3^{J+1})$.
    \item $F_K(v \bmod 3^{K+1})$ is non-empty for all $K \ge 1$ (witnessed by $v \in \D(v)$).
\end{enumerate}
\end{lemma}

\begin{proof}
(1) If $t \in F_K(c)$, $t = (n : \Fp)$ for some $n \in \D(v)$ with $n \equiv c \pmod{3^{K+1}}$. Then $R(n) \in \D(v)$, $R(n) \bmod 3^{K+1} = \rho(c)$ (by reduction), and $(R(n) : \Fp) = \rl(c)(t)$ by Lemma~\ref{lem:operator_compat}(1).

(2) Similar: $L(n) \in \D(v)$, Lemma~\ref{lem:L_cong} gives $L(n) \equiv \lambda(e) \pmod{3^{K+1}}$, and Lemma~\ref{lem:operator_compat}(2) matches the label.

(3) Follows from $3^{J+1} \mid 3^{K+1}$. (4) $v \in \D(v)$ gives $(v : \Fp) \in F_K(v \bmod 3^{K+1})$.
\end{proof}

\begin{proposition}[Cardinality Stabilization]\label{prop:stabilization}
There exist $K_0 \ge 1$ and $\sigma \in \N$ with $\sigma \ge 1$ such that for all $K \ge K_0$ and all valid classes $c < 3^{K+1}$:
\[
|F_K(c)| = \sigma.
\]
\end{proposition}

\begin{proof}
Since right-edge labels have non-zero slopes $4^{\step(c)} \in \Fp^\times$, they act injectively. By Lemma~\ref{lem:fibre_dynamics}(1), $|F_K(c)| \le |F_K(\rho(c))|$. Applying transitivity (Proposition~\ref{prop:transitivity}) in both directions yields:
\[
|F_K(c)| = |F_K(d)|
\]
for all valid classes $c, d < 3^{K+1}$. In particular, $|F_K(c)| = |F_K(8)|$.

Next, by Lemma~\ref{lem:fibre_dynamics}(3) at $J = K$ and class $8$, $F_{K+1}(8) \subseteq F_K(8)$, so the sequence of natural numbers $k \mapsto |F_k(8)|$ is non-increasing. By Lemma~\ref{lem:fibre_dynamics}(4) and transitivity, $|F_K(8)| = |F_K(v \bmod 3^{K+1})| \ge 1$ for all $K$. A non-increasing sequence of positive integers must stabilize: there exist $K_0 \ge 1$ and $\sigma \ge 1$ such that $|F_K(8)| = \sigma$ for all $K \ge K_0$.
\end{proof}

\begin{lemma}[Fibre Equality Across Successive Depths]\label{lem:fibre_succ_eq}
Let $K \ge K_0$. For any valid $e < 3^{K+2}$:
\[
F_{K+1}(e) = F_K(e \bmod 3^{K+1}).
\]
\end{lemma}

\begin{proof}
By Lemma~\ref{lem:fibre_dynamics}(3), $F_{K+1}(e) \subseteq F_K(e \bmod 3^{K+1})$. By Proposition~\ref{prop:stabilization}, both sets are finite and have identical cardinality $\sigma$. A subset of a finite set with equal cardinality is the entire set.
\end{proof}

\needspace{6\baselineskip}
\begin{proposition}[Exact Label Action]\label{prop:walk_image}
Let $K \ge K_0$. For any edge $u \to w$ with label $f$, $f(F_K(u)) = F_K(w)$. Consequently, for any walk from $u$ to $w$ with label $g$:
\[
g(F_K(u)) = F_K(w).
\]
In particular, every holonomy element $h \in \Hol(K)$ acts as a bijection on $F_K(8)$.
\end{proposition}

\begin{proof}
For a right edge $c \to \rho(c)$, $\rl(c)$ is injective and maps $F_K(c)$ into $F_K(\rho(c))$. Since $|F_K(c)| = \sigma = |F_K(\rho(c))|$, the image is the entire target. For a left edge $(e \bmod 3^{K+1}) \to \lambda(e)$, Lemma~\ref{lem:fibre_succ_eq} gives $F_{K+1}(e) = F_K(e \bmod 3^{K+1})$. Thus $\llabel(e)$ maps $F_K(e \bmod 3^{K+1})$ into $F_K(\lambda(e))$, which again is surjective by cardinality comparison. The result for general walks follows by induction on the walk length.
\end{proof}

\section{The Commutator Engine and Simultaneous Density}\label{sec:density}

\begin{lemma}[Commutator Identity]\label{lem:comm_label}
Let $\lambda, \alpha \in \Fp^\times$ and $\beta \in \Fp$. Then:
\[
(\alpha, \beta)^{-1} \then (\lambda, 0)^{-1} \then (\alpha, \beta) \then (\lambda, 0) = \bigl(1, \; (\lambda - 1)\beta\bigr).
\]
\end{lemma}

\begin{proof}
We evaluate the composition stepwise according to Definition~\ref{def:affine_group}:
\begin{align*}
(\alpha, \beta)^{-1} &= (\alpha^{-1}, -\alpha^{-1}\beta), \\
(\lambda, 0)^{-1} &= (\lambda^{-1}, 0).
\end{align*}
Composing the first two factors:
\[
(\alpha^{-1}, -\alpha^{-1}\beta) \then (\lambda^{-1}, 0) = (\lambda^{-1}\alpha^{-1}, -\lambda^{-1}\alpha^{-1}\beta).
\]
Composing with $(\alpha, \beta)$:
\begin{align*}
(\lambda^{-1}\alpha^{-1}, -\lambda^{-1}\alpha^{-1}\beta) \then (\alpha, \beta) 
&= (\alpha \lambda^{-1}\alpha^{-1}, \; \alpha(-\lambda^{-1}\alpha^{-1}\beta) + \beta) \\
&= (\lambda^{-1}, \; (1 - \lambda^{-1})\beta).
\end{align*}
Finally, composing with $(\lambda, 0)$:
\[
(\lambda^{-1}, \; (1 - \lambda^{-1})\beta) \then (\lambda, 0) = (\lambda \lambda^{-1}, \; \lambda(1 - \lambda^{-1})\beta + 0) = (1, \; (\lambda - 1)\beta). \qedhere
\]
\end{proof}

\begin{lemma}[Walk with Non-Zero Translation]\label{lem:ref_walk}
For every $K \ge 1$, there exist $\alpha, \beta \in \Fp$ with $\alpha \neq 0$ and $\beta \neq 0$ such that $(\alpha, \beta) \in \Hol(K)$.
\end{lemma}

\begin{proof}
Consider $e = 8$. Since $8 < 3^{K+2}$, there is a left edge from $8 \bmod 3^{K+1} = 8$ to $\lambda(8) = L(8) \bmod 3^{K+1}$ with label $\llabel(8)$. By Proposition~\ref{prop:transitivity}, choose a right-path from $\lambda(8)$ back to $8$; by Lemma~\ref{lem:walk_of_rIter}, this path has label $(4^E, 0)$ for some $E \in \N$. Composing these yields a closed walk at $8$ with label:
\[
(\alpha, \beta) = \llabel(8) \then (4^E, 0) = \left(4^E \frac{\mult(8)}{3}, \; -4^E \frac{\mult(8)}{6}\right).
\]
Since $p \ge 5$, $4 \neq 0$, and $\mult(8) = 8 \neq 0$, the second coordinate is $\beta = -4^E (8/6) \neq 0$.
\end{proof}

\needspace{8\baselineskip}
\begin{lemma}[Pure Dilation from Loop Pairs]\label{lem:pure_dilation_pair}
Let $e_1, e_2 < 3^{K+2}$ be valid classes. Suppose there exist right-paths:
\begin{itemize}
    \item $u_1 \to u_2$ with $4^{E(n_s, u_1)} = 1$ in $\Fp$, where $u_i = e_i \bmod 3^{K+1}$,
    \item $w_1 \to w_2$ where $w_i = \lambda(e_i)$, with cumulative exponent $E_t = E(n_t, w_1)$.
\end{itemize}
Then:
\[
\left(\frac{\mult(e_2)}{\mult(e_1) \cdot 4^{E_t}}, \; 0\right) \in \Hol(K).
\]
\end{lemma}

\begin{proof}
Let $W_s$ be the walk $u_1 \to u_2$ with label $(1, 0)$, and $W_t$ the walk $w_1 \to w_2$ with label $(4^{E_t}, 0)$. Let $E_1$ be the left edge $u_1 \to w_1$ (label $\llabel(e_1)$) and $E_2$ the left edge $u_2 \to w_2$ (label $\llabel(e_2)$). Form the closed loop based at $w_1$:
\[
\Omega = E_1^{-1} \then W_s \then E_2 \then W_t^{-1}.
\]
Evaluating the label of $\Omega$:
\[
\llabel(e_1)^{-1} = \left(\frac{3}{\mult(e_1)}, \; \frac{1}{2}\right).
\]
Composing with $W_s$ (which acts as the identity $(1, 0)$) leaves the label unchanged. Composing with $\llabel(e_2) = (\mult(e_2)/3, -\mult(e_2)/6)$:
\[
\left(\frac{3}{\mult(e_1)}, \; \frac{1}{2}\right) \then \left(\frac{\mult(e_2)}{3}, \; -\frac{\mult(e_2)}{6}\right) = \left(\frac{\mult(e_2)}{\mult(e_1)}, \; \frac{\mult(e_2)}{3}\left(\frac{1}{2}\right) - \frac{\mult(e_2)}{6}\right) = \left(\frac{\mult(e_2)}{\mult(e_1)}, \; 0\right).
\]
Finally, composing with $W_t^{-1} = (4^{-E_t}, 0)$ gives the label $(\mult(e_2)/(\mult(e_1) 4^{E_t}), 0)$. Connecting $8$ to $w_1$ by a right-path $\gamma$ and conjugating $\Omega$ by $\gamma$ yields a closed walk based at $8$ with identical label.
\end{proof}

\begin{lemma}[Existence of Non-Trivial Dilations]\label{lem:dilation_exists}
For every $K \ge 1$, there exists $\lambda \in \Fp^\times$ with $\lambda \neq 1$ such that $(\lambda, 0) \in \Hol(K)$.
\end{lemma}

\begin{proof}
We divide the analysis into three exhaustive cases:

\textbf{Case 1: $4^{3^K} \neq 1$ in $\Fp$.} By Proposition~\ref{prop:loop8}, there is a closed walk at $8$ with cumulative exponent $E = 3^K$. By Lemma~\ref{lem:walk_of_rIter}, its label is $(4^{3^K}, 0) \in \Hol(K)$. Take $\lambda = 4^{3^K} \neq 1$.

\textbf{Case 2: $4^{3^K} = 1$ in $\Fp$ and $\ord(4) = 3^a$ with $2 \le a \le K$.} Write $a = j + 2$. Set:
\[
e_1 = 8, \qquad e_2 = 8 + 27 \cdot 3^j = 8 + 3^{a+1}.
\]
Both $e_1, e_2$ are valid and bounded by $3^{K+2}$. By Lemma~\ref{lem:mod27}, $\mult(e_1) = \mult(e_2) = 8$. By Lemma~\ref{lem:L_spec}, $L(e_1) = 20$ and $L(e_2) = 20 + 72 \cdot 3^j$. Let $u_i = e_i \bmod 3^{K+1}$ and $w_i = \lambda(e_i)$.
\begin{itemize}
    \item \emph{Source path:} By Proposition~\ref{prop:transitivity}, take a path $u_1 \to u_2$ with exponent $E_s$. Since $e_2 \equiv e_1 \pmod{3^{a+1}}$, Lemma~\ref{lem:iter_mult} forces $4^{E_s} \equiv 1 \pmod{3^{a+1}}$, so $3^a \mid E_s$ by Lemma~\ref{lem:lte}. Since $\ord(4) = 3^a$, $4^{E_s} = 1$ in $\Fp$.
    \item \emph{Target path:} Let $w_1 \to w_2$ have exponent $E_t$. If $4^{E_t} = 1$ in $\Fp$, then $3^a \mid E_t$, which would force $L(e_2) \equiv L(e_1) \pmod{3^{a+1}}$, meaning $72 \cdot 3^j \equiv 0 \pmod{27 \cdot 3^j}$, or $27 \mid 72$, a contradiction. Thus $4^{E_t} \neq 1$ in $\Fp$.
\end{itemize}
Applying Lemma~\ref{lem:pure_dilation_pair} yields $\lambda = 8 / (8 \cdot 4^{E_t}) = 4^{-E_t} \neq 1$ in $\Fp$.

\textbf{Case 3: $4^{3^K} = 1$ in $\Fp$ and $\ord(4) = 3$.} Take $e_1 = 2$ and $e_2 = 11$. Both are valid and $< 3^{K+2}$. Here $\mult(2) = 4$, $\mult(11) = 16$, $L(2) = 2$, and $L(11) = 56$. Notice that:
\[
u_1 \equiv 2, \quad u_2 \equiv 2, \quad w_1 = 2 \equiv 2, \quad w_2 = 56 \equiv 2 \pmod 9.
\]
Because both endpoints of the source path $u_1 \to u_2$ and target path $w_1 \to w_2$ are $\equiv 2 \pmod 9$, Lemma~\ref{lem:iter_mult} reduced modulo $9$ gives $2 \equiv 2 \cdot 4^{E} \pmod 9$ for $E \in \{E_s, E_t\}$, and $4^E \bmod 9$ is $1, 4, 7$ according as $E \bmod 3$ is $0, 1, 2$; hence $3 \mid E_s$ and $3 \mid E_t$. Since $\ord(4) = 3$, $4^{E_s} = 4^{E_t} = 1$ in $\Fp$. Lemma~\ref{lem:pure_dilation_pair} yields:
\[
\lambda = \frac{\mult(11)}{\mult(2) \cdot 1} = \frac{16}{4} = 4 \in \Fp^\times.
\]
Here $\lambda \neq 0$ because $4 \neq 0$ in $\Fp$. If $\lambda = 1$, then $16 = 4$, so $12 = 0$ in $\Fp$, i.e.\ $3 \cdot 4 = 0$. Since $p \ge 5$, both $3 \neq 0$ and $4 \neq 0$, contradicting the absence of zero divisors in $\Fp$. Hence $\lambda \neq 1$.

The cases are exhaustive: either $4^{3^K} \neq 1$ in $\Fp$ (Case~1), or $4^{3^K} = 1$, which forces $\ord(4) \mid 3^K$. In the latter regime, since $p \ge 5$ precludes $\ord(4) = 1$ (as $4 \equiv 1 \pmod p \implies p \mid 3$), we must have $\ord(4) = 3^a$ for some $1 \le a \le K$, splitting into $a = 1$ (Case~3) and $2 \le a \le K$ (Case~2).
\end{proof}

\begin{proposition}[Existence of Translations in Holonomy]\label{prop:translation_exists}
For every $K \ge 1$, there exists $d \in \Fp$ with $d \neq 0$ such that $(1, d) \in \Hol(K)$.
\end{proposition}
\begin{proof}
Choose $(\alpha, \beta) \in \Hol(K)$ with $\alpha, \beta \neq 0$ by Lemma~\ref{lem:ref_walk}. Choose $(\lambda, 0) \in \Hol(K)$ with $\lambda \neq 0, 1$ by Lemma~\ref{lem:dilation_exists}. Since $\Hol(K)$ is a group, it contains the commutator:
\[
(\alpha, \beta)^{-1} \then (\lambda, 0)^{-1} \then (\alpha, \beta) \then (\lambda, 0).
\]
By Lemma~\ref{lem:comm_label}, this element equals $(1, d)$ with $d = (\lambda - 1)\beta$. Since $\Fp$ is a field and $\lambda \neq 1$, $\beta \neq 0$, we have $d \neq 0$.
\end{proof}


\begin{theorem}[Simultaneous Density]\label{thm:density}
Let $p \ge 5$ and $v \in G$. For any $J \ge 1$, any valid class $c < 3^{J+1}$, and any residue $r \in \Fp$, there exists $y \in \D(v)$ such that:
\[
y \equiv c \pmod{3^{J+1}} \quad \text{and} \quad (y : \Fp) = r.
\]
\end{theorem}


\begin{proof}
Let $K_0, \sigma$ be the stabilization parameters from Proposition~\ref{prop:stabilization}. Fix $K = \max(K_0, J)$.


By Proposition~\ref{prop:translation_exists}, there exists a non-trivial translation $(1, d) \in \Hol(K)$ with $d \neq 0$. By Proposition~\ref{prop:walk_image}, $(1, d)$ maps $F_K(8)$ onto itself:
\[
t \in F_K(8) \implies t + d \in F_K(8).
\]
Since $F_K(8)$ is non-empty, choose $t_0 \in F_K(8)$. By induction, $t_0 + k d \in F_K(8)$ for all $k \in \N$. Because $d \in \Fp^\times$, the additive orbit $\{t_0 + k d : k \in \N\}$ equals all of $\Fp$. Thus $F_K(8) = \Fp$.


This implies that the stabilized cardinality is $\sigma = |F_K(8)| = p$. Since $c < 3^{J+1} \le 3^{K+1}$, $c$ is a valid class at depth $K$, so $|F_K(c)| = \sigma = p$, which forces $F_K(c) = \Fp$. Thus $r \in F_K(c)$.


Since $J \le K$, Lemma~\ref{lem:fibre_dynamics}(3) gives $r \in F_J(c \bmod 3^{J+1}) = F_J(c)$. By Definition~\ref{def:fibres}, there exists $y \in \D(v)$ satisfying $y \equiv c \pmod{3^{J+1}}$ and $(y : \Fp) = r$.
\end{proof}


\section{Small Primes and the Master Theorem}\label{sec:small_primes}


\begin{proposition}[Case $p \ge 5$]\label{prop:p_ge_5}
Theorem~\ref{thm:main} holds for $p \ge 5$.
\end{proposition}


\begin{proof}
By Lemma~\ref{lem:seed_exist}, choose $v \in G$ reaching $w$. Let $J = B + 1$ and define $c = 9 \cdot 3^B - 1 = 3^{B+2} - 1$. Then $c \equiv 8 \pmod 9$, so $c$ is valid; $c < 3^{J+1}$; and $c > B$ (since $B < 3^B \le c$). Applying Theorem~\ref{thm:density} with class $c$ and residue $r \in \Fp$ yields $y \in \D(v)$ such that $y \equiv c \pmod{3^{J+1}}$ and $(y : \Fp) = r$.


Since $y \bmod 3^{J+1} = c$ and $c > 0$, we have $y \ge c > B$. By Lemma~\ref{lem:desc_sound}, $y \reaches v$, and since $v \reaches w$, transitivity ensures $y \reaches w$.
\end{proof}


\begin{proposition}[Case $p = 2$]\label{prop:p_eq_2}
Theorem~\ref{thm:main} holds for $p = 2$.
\end{proposition}


\begin{proof}
Let $r \in \Z/2\Z$.
\begin{itemize}[noitemsep, topsep=1pt]
    \item If $r = 0$: set $n = 2^{B+1} w$. Then $n > B$, $n \equiv 0 \pmod 2$, and $n \reaches w$ in $B+1$ halvings.
    \item If $r = 1$: choose $t = 3B + 3$. There exists $k \in \{2t, 2t+1\}$ such that $2^k w \equiv 1 \pmod 3$ (if $w \equiv 1 \pmod 3$, choose $k = 2t$; if $w \equiv 2 \pmod 3$, choose $k = 2t+1$). Since $k \ge 6B + 6$, $2^k w > 6B + 6$. Define $n = (2^k w - 1)/3$. Then $3n + 1 = 2^k w$. Since $2^k w$ is even, $n$ is odd, so $(n : \Z/2\Z) = 1$. Furthermore, $n > B$, and $\C(n) = 2^k w$, which reaches $w$ in $k$ halvings. \qedhere
\end{itemize}
\end{proof}


\needspace{6\baselineskip}
\begin{proposition}[Case $p = 3$]\label{prop:p_eq_3}
Theorem~\ref{thm:main} holds for $p = 3$.
\end{proposition}


\begin{proof}
Let $r \in \Z/3\Z$.
\begin{itemize}[noitemsep, topsep=1pt]
    \item If $r = 0$: $2$ is a primitive root modulo $9$. There exists $j < 6$ such that $2^j w \equiv 1 \pmod 9$. Set $k = j + 6(B+1)$. Since $64 \equiv 1 \pmod 9$, $2^k w = 64^{B+1} \cdot 2^j w \equiv 1 \pmod 9$. Set $n = (2^k w - 1)/3$. Since $2^k w$ is even, $n$ is odd, so $\C(n) = 3n + 1 = 2^k w$, and $n \reaches w$. Since $2^k w \equiv 1 \pmod 9$, $n \equiv 0 \pmod 3$. Finally, $n > B$.
    \item If $r = 1$: choose $k \ge B$ such that $2^k w \equiv 1 \pmod 3$. Set $n = 2^k w$. Then $n > B$, $n \equiv 1 \pmod 3$, and $n \reaches w$ in $k$ halvings.
    \item If $r = 2$: with $k$ as above, set $n = 2^{k+1} w$. Then $n \equiv 2 \cdot 1 = 2 \pmod 3$, $n > B$, and $n \reaches w$. \qedhere
\end{itemize}
\end{proof}


\begin{proof}[Proof of Theorem~\ref{thm:main}]
Combine Propositions~\ref{prop:p_ge_5}, \ref{prop:p_eq_2}, and \ref{prop:p_eq_3}.
\end{proof}


\needspace{12\baselineskip}
\section{Application: Local Obstruction Barriers for Cycles}\label{sec:cycles}


A hypothetical Collatz cycle with $k$ odd steps and $S$ total steps satisfies the Diophantine equation:
\[
(2^S - 3^k) n = N,
\]
where $N = \sum_{i=0}^{k-1} 3^{k-1-i} 2^{s_i}$ is a positive integer determined by the step sequence.


\begin{definition}[Cycle Specification]\label{def:cycle_spec}
A \emph{cycle specification} is a triple $(k, S, N) \in \N^3$.
\end{definition}


\begin{theorem}[Predecessors Modulo $p$]\label{thm:reaching_dense}
Let $p$ be any prime. If $v > 0$ with $3 \nmid v$, then for every $r \in \Fp$, there exists $n \in \N$ such that $n \reaches v$ and $(n : \Fp) = r$.
\end{theorem}


\begin{proof}
Specialization of Theorem~\ref{thm:main} with threshold $B = 0$.
\end{proof}


\begin{theorem}[Absence of Local Obstructions]\label{thm:barrier}
Let $p$ be any prime. Let $(k, S, N)$ be a cycle specification such that
\[
(2^S - 3^k : \Fp) \neq 0.
\]
Then for every target $v > 0$ with $3 \nmid v$, there exists $n \in \N$ such that $n \reaches v$ and
\[
(2^S - 3^k : \Fp) \cdot (n : \Fp) = (N : \Fp).
\]
\end{theorem}


\begin{proof}
Since $(2^S - 3^k : \Fp) \neq 0$ and $\Fp$ is a field, the linear equation has a unique solution:
\[
r = (2^S - 3^k : \Fp)^{-1} (N : \Fp) \in \Fp.
\]
By Theorem~\ref{thm:reaching_dense}, there exists $n \in \N$ reaching $v$ with $(n : \Fp) = r$. Multiplying by $(2^S - 3^k : \Fp)$ gives $(2^S - 3^k : \Fp) \cdot (n : \Fp) = (N : \Fp)$.
\end{proof}


\begin{remark}\label{rem:barrier}
Theorem~\ref{thm:barrier} demonstrates that for any prime $p$ where the cycle denominator does not vanish, modular conditions modulo $p$ alone cannot rule out the existence of cycle predecessors. The global obstruction to non-trivial cycles must necessarily involve multi-prime interactions or archimedean constraints.
\end{remark}


% ---------------------------------------------------------------------------
% APPENDIX
% ---------------------------------------------------------------------------
\clearpage
\appendix
\section*{Formalization Correspondence Table}

{\footnotesize
\setlength{\tabcolsep}{3pt}
\renewcommand{\arraystretch}{0.95}
\begin{longtable}{@{}>{\raggedright\arraybackslash}p{0.44\textwidth}>{\raggedright\arraybackslash}p{0.52\textwidth}@{}}
\caption{Correspondence between paper mathematical statements and verified declarations in \texttt{ERE.lean}.}\label{tab:correspondence}\\
\toprule
\textbf{Paper Definition / Result} & \textbf{Lean 4 Formalization (\texttt{ERE.*})} \\
\midrule
\endfirsthead
\toprule
\textbf{Paper Definition / Result} & \textbf{Lean 4 Formalization (\texttt{ERE.*})} \\
\midrule
\endhead
\bottomrule
\endfoot
\bottomrule
\endlastfoot
Collatz operator $\C$, reachability $\reaches$ & \texttt{collatz}, \texttt{Reaches} \\
Lemma~\ref{lem:atomic} (Halving, atomic inverse steps) & \texttt{collatz\_double}, \texttt{collatz\_pow\_mul}, \texttt{reaches\_trans}, \texttt{collatz\_eq\_iff} \\
Corollary~\ref{cor:three_dead} (Multiples of $3$) & \texttt{pred\_of\_three\_dvd} \\
Good numbers, coordinate $x(n)$ & \texttt{GoodOdd}, \texttt{toX} \\
Lemma~\ref{lem:coord_step} (Coordinate step) & \texttt{collatz\_odd}, \texttt{reaches\_toX} \\
Valid class, domain $G$ & \texttt{ValidClass}, \texttt{InG} \\
Lemma~\ref{lem:good_valid} (Good numbers are valid classes) & \texttt{toX\_inG}, \texttt{exists\_goodOdd\_of\_inG}, \texttt{inG\_iff} \\
Branches $R, L$, multiplier $\mult$ & \texttt{R}, \texttt{L}, \texttt{mult} \\
Abstract descendant tree $\D(v)$ & \texttt{Desc} \\
Proposition~\ref{prop:branch_derivation} (Derivation of branches) & \texttt{R\_first}, \texttt{L\_first}, \texttt{collatz\_third\_pred}, \texttt{Rorig}, \texttt{multOrig}, \texttt{Lorig}, \texttt{R\_conj}, \texttt{L\_conj} \\
Lemma~\ref{lem:collatz_elementary} (Branch trajectories) & \texttt{R\_reaches}, \texttt{L\_reaches} \\
Lemma~\ref{lem:desc_sound} (Soundness of $\D$) & \texttt{Desc\_reaches} \\
Lemma~\ref{lem:seed_exist} (Seed existence) & \texttt{exists\_inG\_reaches} \\
Lemma~\ref{lem:L_spec} (Exactness of $L$) & \texttt{L\_spec} \\
Lemma~\ref{lem:mod27} (Residues mod $27$) & \texttt{residue\_mod27}, \texttt{mult\_mem}, \texttt{mult\_pow} \\
Lemma~\ref{lem:G_inv} (Invariance of $G$) & \texttt{R\_mem}, \texttt{L\_mem}, \texttt{L\_valid}, \texttt{Desc\_mem} \\
Lemma~\ref{lem:L_cong} (Precision loss under $L$) & \texttt{mod27\_eq\_of\_mod\_pow}, \texttt{mult\_congr}, \texttt{L\_cong} \\
$\rho, \lambda, \step, \rho^n, E(n, c)$ & \texttt{rnext}, \texttt{lnext}, \texttt{rStep}, \texttt{rIter}, \texttt{rExp} \\
Lemma~\ref{lem:iter_mult} (Multiplication action) & \texttt{rIter\_eq}, \texttt{\detokenize{rIter_eq'}} \\
Lemma~\ref{lem:step_pair} (Two-step tripling) & \texttt{rStep\_pair}, \texttt{rExp\_two\_mul} \\
Lemma~\ref{lem:lte} (LTE for base $4$) & \texttt{lte\_four}, \texttt{four\_pow\_modEq\_one\_iff} \\
Lemma~\ref{lem:four_cancel} (Four-cancellation) & \texttt{four\_cancel} \\
Lemma~\ref{lem:exists_exp} (Exponent classification) & \texttt{four\_pow\_mod9}, \texttt{exp\_type}, \texttt{exists\_exp} \\
Proposition~\ref{prop:transitivity} (Transitivity of $\rho$) & \texttt{R\_transitive} \\
Proposition~\ref{prop:loop8} (Loop at $8$) & \texttt{R\_full\_loop} \\
Affine group $\Aff(p)$, holonomy $\Hol(K)$ & \texttt{Aff}, \texttt{aapp}, \texttt{athen}, \texttt{ainv}, \texttt{Edge}, \texttt{Walk}, \texttt{Hol} \\
Lemma~\ref{lem:operator_compat} (Operator equivariance) & \texttt{R\_label}, \texttt{L\_label}, \texttt{rightLabel\_fst}, \texttt{rightLabel\_snd} \\
Lemma~\ref{lem:walk_of_rIter} (Walks from right iterates) & \texttt{walk\_of\_rIter} \\
Fibres $F_K(c)$ & \texttt{M}, \texttt{fibre} \\
Lemma~\ref{lem:fibre_dynamics} (Fibre dynamics) & \texttt{fibre\_right}, \texttt{fibre\_left}, \texttt{fibre\_proj}, \texttt{fibre\_root\_nonempty} \\
Proposition~\ref{prop:stabilization} (Fibre stabilization) & \texttt{fibre\_ncard\_step}, \texttt{fibre\_ncard\_chain}, \texttt{fibre\_ncard\_const}, \texttt{stabilization} \\
Lemma~\ref{lem:fibre_succ_eq} (Successive depths) & \texttt{fibre\_succ\_eq} \\
Proposition~\ref{prop:walk_image} (Exact label action) & \texttt{edge\_image}, \texttt{walk\_image} \\
Lemma~\ref{lem:comm_label} (Commutator identity) & \texttt{comm\_label} \\
Lemma~\ref{lem:ref_walk} (Reference walk) & \texttt{ref\_walk} \\
Lemma~\ref{lem:pure_dilation_pair} (Pair dilation) & \texttt{rpath}, \texttt{pure\_dilation\_of\_pair} \\
Lemma~\ref{lem:dilation_exists} (Dilation existence) & \texttt{dil\_case1}, \texttt{dil\_case2a}, \texttt{dil\_case2b}, \texttt{order\_split}, \texttt{dilation\_exists} \\
Proposition~\ref{prop:translation_exists} (Translation in $\Hol$) & \texttt{translation\_of\_dilation}, \texttt{translation\_exists} \\
Theorem~\ref{thm:density} (Simultaneous density) & \texttt{main\_theorem\_core} \\
Proposition~\ref{prop:p_ge_5} ($p \ge 5$) & \texttt{main\_theorem\_p\_ge\_5} \\
Proposition~\ref{prop:p_eq_2} ($p = 2$) & \texttt{main\_theorem\_p\_eq\_2} \\
Proposition~\ref{prop:p_eq_3} ($p = 3$) & \texttt{main\_theorem\_p\_eq\_3} \\
Theorem~\ref{thm:main} (Master Theorem) & \texttt{main\_theorem\_infinite} \\
Cycle Specification & \texttt{CycleSpec} \\
Theorem~\ref{thm:reaching_dense} (Reaching density) & \texttt{reaching\_predecessor\_dense\_mod\_p} \\
Theorem~\ref{thm:barrier} (Cycle barrier) & \texttt{cycle\_equation\_solvable\_mod\_p}, \texttt{cycle\_local\_obstruction\_barrier} \\
\end{longtable}
}

{\footnotesize
\setlength{\tabcolsep}{4.5pt}
\renewcommand{\arraystretch}{1.08}
\begin{longtable}{@{}>{\raggedright\arraybackslash}p{0.26\textwidth}>{\raggedright\arraybackslash}p{0.69\textwidth}@{}}
\caption{Auxiliary Lean~4 declarations supporting the primary equivalence.}\label{tab:auxiliary}\\
\toprule
\textbf{Auxiliary Facts} & \textbf{Lean 4 Declarations (\texttt{ERE.*})} \\
\midrule
\endfirsthead
\toprule
\textbf{Auxiliary Facts} & \textbf{Lean 4 Declarations (\texttt{ERE.*})} \\
\midrule
\endhead
\bottomrule
\endfoot
\bottomrule
\endlastfoot
Classes, exponents & \texttt{nine\_dvd\_pow}, \texttt{validClass\_mod}, \texttt{valid8}, \texttt{eight\_lt}, \texttt{rnext\_valid}, \texttt{rIter\_valid}, \texttt{rIter\_modEq}, \texttt{exp\_dvd\_of\_congr}, \texttt{congr\_of\_exp\_dvd}, \texttt{rExp\_mod3}, \texttt{four\_pow\_eq\_one\_of\_mod3} \\
Affine labels, walks, fibres & \texttt{athen\_app}, \texttt{athen\_assoc}, \texttt{athen\_one\_left}, \texttt{athen\_one\_right}, \texttt{ainv\_app}, \texttt{ainv\_ainv}, \texttt{ainv\_athen}, \texttt{aapp\_inj}, \texttt{two\_ne\_zero\_zmod}, \texttt{three\_ne\_zero\_zmod}, \texttt{four\_ne\_zero\_zmod}, \texttt{six\_ne\_zero\_zmod}, \texttt{mult\_cast\_ne\_zero}, \texttt{Edge.label\_unit}, \texttt{Walk.label\_unit}, \texttt{Walk.append}, \texttt{Walk.symm}, \texttt{Hol\_mul}, \texttt{Hol\_inv}, \texttt{HasTranslation}, \texttt{R\_class}, \texttt{L\_class}, \texttt{rightLabel\_mod}, \texttt{leftLabel\_mod}, \texttt{fibre\_ncard\_antitone} \\
Collatz bridge, small primes & \texttt{collatz\_8}, \texttt{inverse\_tree\_covers}, \texttt{exists\_large\_valid\_class}, \texttt{cast\_eq\_of\_mod}, \texttt{\detokenize{cast_eq_of_mod'}}, \texttt{exists\_pow\_two\_mod}, \texttt{exists\_k\_mod\_three}, \texttt{four\_pow\_mod\_three}, \texttt{sixtyfour\_pow\_mod}, \texttt{lt\_pow\_mul}, \texttt{odd\_of\_three\_mul\_add\_one}, \texttt{reaches\_of\_three\_mul\_add\_one}, \texttt{ge\_of\_mod\_eq} \\
\end{longtable}
}

\subsection*{Formalization Architecture and Verification Environment}
{\emergencystretch=3em The complete formalization is self-contained in \texttt{ERE.lean}. The development introduces zero custom axioms and contains no unproven assertions (\texttt{sorry}). Running Lean's \texttt{\#print axioms main\_theorem\_infinite} confirms that the verification rests strictly upon standard classical foundations: \texttt{propext}, \texttt{Classical.choice}, and \texttt{Quot.sound}. The project builds under Lean~4 toolchain \texttt{leanprover/lean4:v4.36.0-rc1}, Mathlib commit \texttt{0b57f69}, Batteries \texttt{8e746d7} and Aesop \texttt{29c64c5}. The repository, including \texttt{lakefile.toml} and \texttt{lake-manifest.json}, is available at\par}
\begin{center}
\url{https://github.com/Sio-MLST/Every-Remainder-Everywhere}
\end{center}

\end{document}
